\section{Penentuan Bye pada Bracket Eliminasi}

Dalam kompetisi panahan yang menggunakan sistem eliminasi (seperti pada IANSEO atau World Archery), jumlah peserta tidak selalu sama dengan ukuran bracket eliminasi. Oleh karena itu, digunakan mekanisme yang disebut \textbf{bye}.

\subsection{Ukuran Bracket}

Bracket eliminasi biasanya menggunakan ukuran berbasis pangkat dua:

\[
4, 8, 16, 32, 64, 128
\]

Jika jumlah peserta tidak sama dengan ukuran tersebut, maka dipilih ukuran bracket terkecil yang lebih besar atau sama dengan jumlah peserta.

\textbf{Contoh:}

\begin{itemize}
\item Jumlah peserta = 60
\item Ukuran bracket terdekat = 64
\end{itemize}

\subsection{Perhitungan Bye}

Jumlah bye dihitung dengan rumus:

\[
\text{Bye} = \text{Bracket Size} - \text{Jumlah Peserta}
\]

Contoh:

\[
64 - 60 = 4
\]

Sehingga terdapat \textbf{4 bye} pada ronde pertama.

\subsection{Distribusi Bye}

Bye diberikan kepada peserta dengan \textbf{ranking tertinggi} berdasarkan hasil qualification. Hal ini dilakukan untuk menjaga keadilan kompetisi.

Jika terdapat 4 bye, maka:

\begin{itemize}
\item Seed 1 mendapat bye
\item Seed 2 mendapat bye
\item Seed 3 mendapat bye
\item Seed 4 mendapat bye
\end{itemize}

Peserta tersebut langsung maju ke ronde berikutnya tanpa bertanding pada ronde pertama.

\subsection{Pola Pairing Eliminasi}

Setelah bye diberikan, pertandingan lainnya mengikuti pola pairing standar:

\[
1 \text{ vs } N, \quad 2 \text{ vs } (N-1), \quad 3 \text{ vs } (N-2)
\]

Dengan contoh bracket 64:

\begin{verbatim}
Seed 1  vs Bye
Seed 2  vs Bye
Seed 3  vs Bye
Seed 4  vs Bye
Seed 5  vs Seed 60
Seed 6  vs Seed 59
Seed 7  vs Seed 58
...
\end{verbatim}

\subsection{Alur Pembuatan Bracket Eliminasi}

Secara umum proses pembentukan bracket eliminasi adalah sebagai berikut:

\begin{verbatim}
Qualification Ranking
        ↓
Menentukan ukuran bracket (power of 2)
        ↓
Menghitung jumlah bye
        ↓
Memberikan bye ke seed tertinggi
        ↓
Membentuk pasangan pertandingan eliminasi
\end{verbatim}

\subsection{Kesimpulan}

\begin{itemize}
\item Ukuran bracket eliminasi mengikuti bilangan pangkat dua.
\item Jika jumlah peserta lebih sedikit dari ukuran bracket, maka digunakan bye.
\item Bye diberikan kepada peserta dengan ranking qualification tertinggi.
\item Pairing pertandingan mengikuti pola seed untuk menjaga fairness kompetisi.
\end{itemize}